\cleardoublepage
\vspace{2cm}
\begin{center}
{\fontsize{14}{16.8}\selectfont\textbf{DAFTAR SIMBOL}}\\[1em]
\end{center}
\label{chap:simbol}
\addcontentsline{toc}{chapter}{DAFTAR SIMBOL}
\setcounter{chapter}{0}
\setcounter{section}{0}
\vspace{2em}

% Simbol pada Persamaan V.1--V.4 (Bab V); deskripsi mengikuti teks docx.
% Kolom "Pemakaian pertama kali" = halaman persamaan tempat simbol pertama muncul.
\begin{longtable}{@{} l >{\raggedright\arraybackslash}p{8cm} c @{}}
\multicolumn{1}{l}{\textbf{Notasi}} & \multicolumn{1}{l}{\textbf{Deskripsi}} & \multicolumn{1}{l}{\parbox[c]{3.5cm}{\raggedright \textbf{Pemakaian pertama kali}}} \\[1ex]
\endfirsthead
\multicolumn{1}{l}{\textbf{Notasi}} & \multicolumn{1}{l}{\textbf{Deskripsi}} & \multicolumn{1}{l}{\parbox[c]{3.5cm}{\raggedright \textbf{Pemakaian pertama kali}}} \\[1ex]
\endhead
$c$ & Jumlah kategori penilaian & 148 \\
$\mathit{lo}$ & Skor terendah & 148 \\
$n$ & Jumlah butir pada level abstraksi yang dinilai (Persamaan~\ref{eq:V.1}); jumlah pakar (Persamaan~\ref{eq:V.3}) & 146 \\
$r$ & Skor dari pakar & 148 \\
$s$ & $r - \mathit{lo}$ & 148 \\
$V$ & Indeks Aiken’s V & 148 \\
\end{longtable}
